\clearpage
\VerdierPhysicalPageAnchor{0144}%
\begingroup
\small
\setlength{\parskip}{0pt}%
\setlength{\emergencystretch}{1.5em}%
\setlength{\abovedisplayskip}{3pt plus 1pt minus 1pt}%
\setlength{\belowdisplayskip}{3pt plus 1pt minus 1pt}%

\VerdierSemanticAnchor{ch2-2-3-1-proof-c-bis-converse}
\noindent saturée. La sous-catégorie $Q(Q^{-1}(\mathcal C))$ est strictement
pleine saturée d'après la démonstration de d). Or la catégorie
$Q(Q^{-1}(\mathcal C))$ est la clôture strictement pleine de $\mathcal C$.
Par suite, $\mathcal C$ est saturée. Ceci démontre l'assertion
c)$^{\mathrm{bis}}$.

\VerdierSemanticAnchor{ch2-2-3-1-proof-d-bis}
Il reste à démontrer la dernière assertion. Il est clair que lorsque
$\mathcal C$ est strictement pleine, on a $Q(Q^{-1}(\mathcal C))=\mathcal C$.
Soit $\mathcal A$ une catégorie triangulée strictement pleine contenant
$\widehat{\mathcal B}$. Il nous faut montrer que
$\mathcal A=Q^{-1}(Q(\mathcal A))$. Quitte à remplacer $\mathcal B$ par
$\widehat{\mathcal B}$, ce qui ne change pas la catégorie
$\mathcal D/\mathcal B$ (2.2.11), on peut supposer que $\mathcal B$ est
saturée, et par suite que $\mathcal B=\widehat{\mathcal B}$. Comme on a
évidemment l'inclusion
$\mathcal A\subset Q^{-1}(Q(\mathcal A))$, il suffit de montrer que tout
objet de $\mathcal D$ qui devient isomorphe dans $\mathcal D/\mathcal B$ à
un objet de $\mathcal A$ est un objet de $\mathcal A$.

\VerdierSemanticAnchor{ch2-2-3-1-proof-d-bis-roof-setup}
Soient donc $X$ un objet de $\mathcal A$, $Y$ un objet de $\mathcal D$ et
$m\colon Q(X)\xrightarrow{\sim}Q(Y)$ un isomorphisme. Le morphisme $m$ se
décompose en $m=Q(f)Q(s)^{-1}$, où
$s\in S_{\mathcal D}(\mathcal B)$ (2.2.2), et où $Q(f)$ est un
isomorphisme. Donc $f\in S_{\mathcal D}(\mathcal B)$ car $\mathcal B$ est
saturée ((2.1.8) et (2.2.4)). On a donc un diagramme dans $\mathcal D$ :

\VerdierSemanticAnchor{ch2-2-3-1-proof-d-bis-roof}
\begin{center}
  \begin{tikzcd}[column sep=4.2em,row sep=2.5em,baseline=(current bounding box.center)]
    & Z \arrow[dl,"s"'] \arrow[dr,"f"] & \\
    X && Y
  \end{tikzcd}
  \qquad $s,f\in S_{\mathcal D}(\mathcal B)$\quad .
\end{center}

\VerdierSemanticAnchor{ch2-2-3-1-proof-d-bis-complete}
Le morphisme $s$ étant un élément de $S_{\mathcal D}(\mathcal B)$, il
existe un triangle distingué $(Z,X,W,s,u,v)$, où $W$ est un objet de
$\mathcal B$ (2.1.7). Par suite, $W$ est un objet de $\mathcal A$, et comme
$\mathcal A$ est strictement pleine, l'objet $Z$ appartient aussi à
$\mathcal A$ (1.2.4). Le même argument montre alors que $Y$ est un objet de
$\mathcal A$, ce qui achève la démonstration de la proposition.

\VerdierSemanticAnchor{ch2-2-3-2}
\medskip
\noindent\textbf{2.3.2.}\quad
\VerdierSemanticAnchor{ch2-2-3-2-free-object}
Soient $F\colon C\longrightarrow C'$ un foncteur et $X$ un objet de $C'$.
Rappelons qu'un objet $Y$ de $C$ est dit \emph{$F$-libre à droite sur
$X$}, ou simplement \emph{libre à droite sur $X$}, s'il existe un
morphisme $u\colon F(Y)\longrightarrow X$ tel que pour tout morphisme
$F(Z)\xrightarrow{v}X$, il existe un et un seul morphisme
$w\colon Z\longrightarrow Y$ tel que le diagramme ci-après soit
commutatif :

\VerdierSemanticAnchor{ch2-2-3-2-free-object-diagram}
\begin{center}
  \begin{tikzcd}[column sep=4.7em,row sep=2.7em]
    F(Y) \arrow[rr,"u"] && X \\
    & F(Z) \arrow[ul,"{F(w)}"] \arrow[ur,"v"'] &
  \end{tikzcd}
  \quad .
\end{center}

\VerdierSemanticAnchor{ch2-2-3-2-freedom-arrow-boundary}
La flèche $u\colon F(Y)\longrightarrow X$ s'appelle \emph{flèche de
liberté} pour le couple $(X,Y)$ (relativement au foncteur $F$). Cela peut
encore s'interpréter, en prenant un
\endgroup
